% --- Archon physics formalization source begin ---
% archon:physics
% archon:covers PhyXMiniProblems/problem_phyx_mini_0602.lean
% archon:source-report reports/phyx_mini/problem_phyx_mini_0602.source.json

\chapter{Physics problem phyx\_mini\_0602}
\label{ch:PhyXMiniProblems_problem_phyx_mini_0602}

\paragraph{Problem source.}
\#\# Physical scenario

For some purposes it is more convenient to work with the transfer matrix, \$\textbackslash{}mathbf\{M\}\$, which gives you the amplitudes to the right of the potential (\$F\$ and \$G\$) in terms of those to the left (\$A\$ and \$B\$): \$\$egin\{pmatrix\} F \textbackslash{} G \textbackslash{}end\{pmatrix\} = egin\{pmatrix\} M\_\{11\} \& M\_\{12\} \textbackslash{} M\_\{21\} \& M\_\{22\} \textbackslash{}end\{pmatrix\} egin\{pmatrix\} A \textbackslash{} B \textbackslash{}end\{pmatrix\}.\$\$

\#\# Question

Find the \$M\$-matrix for scattering from the double delta-function \$\$V(x) = -\textbackslash{}alpha [\textbackslash{}delta(x+a) + \textbackslash{}delta(x-a)].\$\$ What is the transmission coefficient for this potential?

\#\# Answer choices

- A: A: \textbackslash{}left[ 1 + \textbackslash{}frac\{m\textbackslash{}alpha\textasciicircum{}2\}\{\textbackslash{}hbar\textasciicircum{}2 E\} \textbackslash{}left( \textbackslash{}sin(2ka) + \textbackslash{}frac\{\textbackslash{}hbar k\}\{m\textbackslash{}alpha\} 
ight)\textasciicircum{}2 
ight]\textasciicircum{}\{-1\}

- B: B: \textbackslash{}left[ 1 + \textbackslash{}left( \textbackslash{}frac\{2\textbackslash{}alpha\}\{k\} 
ight)\textasciicircum{}2 \textbackslash{}cos\textasciicircum{}2(ka) 
ight]\textasciicircum{}\{-1\}

- C: C: \textbackslash{}left[ 1 + \textbackslash{}frac\{2m\textbackslash{}alpha\textasciicircum{}2\}\{\textbackslash{}hbar\textasciicircum{}2 E\} \textbackslash{}left( \textbackslash{}cos(2ka) - \textbackslash{}frac\{m\textbackslash{}alpha\}\{\textbackslash{}hbar k\} \textbackslash{}sin(2ka) 
ight) 
ight]\textasciicircum{}\{-1\}

- D: D: \textbackslash{}left[ \textbackslash{}frac\{1\}\{1 + \textbackslash{}left( \textbackslash{}frac\{m\textbackslash{}alpha\}\{\textbackslash{}hbar\textasciicircum{}2\} 
ight)\textasciicircum{}2\} 
ight]

\#\# Figure caption (auxiliary; use the image as primary evidence)

This image depicts a graphical representation with a horizontal axis labeled \textbackslash{}( x \textbackslash{}), which extends from left to right. The graph features two distinct peaks, labeled as \textbackslash{}( M\_1 \textbackslash{}) and \textbackslash{}( M\_2 \textbackslash{}).

- The first peak (\textbackslash{}( M\_1 \textbackslash{})) is located on the left side, and the curve has a pronounced amplitude, showing undulations.
- The second peak (\textbackslash{}( M\_2 \textbackslash{})) appears on the right side, also with undulations but varying in shape from \textbackslash{}( M\_1 \textbackslash{}).

The graph is divided into three distinct segments:

1. **Left Segment:** 
   - Text: \textbackslash{}( V = 0 \textbackslash{})
   - Location: Below a bracket that spans a flat line section before the first peak.
   
2. **Middle Segment:** 
   - Text: \textbackslash{}( V = 0 \textbackslash{})
   - Location: Under a bracket that spans the region between \textbackslash{}( M\_1 \textbackslash{}) and \textbackslash{}( M\_2 \textbackslash{}), inclusive, of the valleys following each peak.
   
3. **Right Segment:** 
   - Text: \textbackslash{}( V = 0 \textbackslash{})
   - Location: Below a bracket beyond the second peak, indicating another flat line segment.

This arrangement suggests a potential representation of a physical phenomenon, like a potential energy diagram in physics, where \textbackslash{}( V = 0 \textbackslash{}) indicates areas of zero potential energy.

\#\# Dataset metadata

Quantum Phenomena; Numerical Reasoning, Physical Model Grounding Reasoning

\paragraph{Recorded answer/context.}
C: C: \textbackslash{}left[ 1 + \textbackslash{}frac\{2m\textbackslash{}alpha\textasciicircum{}2\}\{\textbackslash{}hbar\textasciicircum{}2 E\} \textbackslash{}left( \textbackslash{}cos(2ka) - \textbackslash{}frac\{m\textbackslash{}alpha\}\{\textbackslash{}hbar k\} \textbackslash{}sin(2ka) 
ight) 
ight]\textasciicircum{}\{-1\}

\paragraph{Figure/image path.}
/root/proposal\_for\_physic/science-mango/phyx\_mini\_run/phyx\_data/test\_image/602.png

\paragraph{Formalization target.}
create a compiling Lean file with sorry bodies at `PhyXMiniProblems/problem\_phyx\_mini\_0602.lean`. The Lean declarations must preserve the physical quantities, dimensions or dimensional roles, figure labels, governing-law hypotheses, and final relation expressed by this problem.
Use Mathlib/Physlib names found through LeanExplore where available. If a physics API is missing, introduce faithful local abstractions rather than scalar placeholder aliases.

\begin{theorem}[Physics formalization target]
\label{thm:physics:phyx_mini_0602:target}
\lean{PhyXMiniProblems.ProblemPhyXMini0602.problem_phyx_mini_0602}
\uses{def:physics:phyx-mini-0602:phyxminiproblems-problemphyxmini0602-doubledeltascatteringsetup, def:physics:phyx-mini-0602:phyxminiproblems-problemphyxmini0602-halfseparationmeters, def:physics:phyx-mini-0602:phyxminiproblems-problemphyxmini0602-matchesdoubledeltascenario, def:physics:phyx-mini-0602:phyxminiproblems-problemphyxmini0602-matchessupplieddoubledeltafigure, def:physics:phyx-mini-0602:phyxminiproblems-problemphyxmini0602-matchesleftincidentscatteringstate, def:physics:phyx-mini-0602:phyxminiproblems-problemphyxmini0602-hasphysicaldoubledeltaparameters, def:physics:phyx-mini-0602:phyxminiproblems-problemphyxmini0602-usesstandardreducedplanckconstant, def:physics:phyx-mini-0602:phyxminiproblems-problemphyxmini0602-satisfiesfreeparticledispersionlaw, def:physics:phyx-mini-0602:phyxminiproblems-problemphyxmini0602-satisfiesdoubledeltascatteringlaws, def:physics:phyx-mini-0602:phyxminiproblems-problemphyxmini0602-attractivedeltatransfermatrix, def:physics:phyx-mini-0602:phyxminiproblems-problemphyxmini0602-explicitdoubledeltatransfermatrix, def:physics:phyx-mini-0602:phyxminiproblems-problemphyxmini0602-physicallysupportedtransmissioncoefficient}
The assigned autoformalize agent should translate this physics problem into Lean declarations in the covered file, with theorem and lemma proof bodies written as `by sorry`.
\end{theorem}
\begin{proof}
This is an autoformalization task, not a proof task. Produce statements that can later be proved by the physics prover without weakening the physical model.
\end{proof}

% --- Archon declaration topology begin ---
\section{Declaration topology}
\begin{definition}[Length Quantity]
  \label{def:physics:phyx-mini-0602:phyxminiproblems-problemphyxmini0602-lengthquantity}
  \lean{PhyXMiniProblems.ProblemPhyXMini0602.LengthQuantity}
  A nonnegative, unit-independent physical length
\end{definition}
\begin{proof}
  This is the declared model definition of Length Quantity.
\end{proof}
\begin{definition}[Mass Quantity]
  \label{def:physics:phyx-mini-0602:phyxminiproblems-problemphyxmini0602-massquantity}
  \lean{PhyXMiniProblems.ProblemPhyXMini0602.MassQuantity}
  A nonnegative, unit-independent particle mass
\end{definition}
\begin{proof}
  This is the declared model definition of Mass Quantity.
\end{proof}
\begin{definition}[action Dimension]
  \label{def:physics:phyx-mini-0602:phyxminiproblems-problemphyxmini0602-actiondimension}
  \lean{PhyXMiniProblems.ProblemPhyXMini0602.actionDimension}
  The physical dimension of action, mass * length² / time
\end{definition}
\begin{proof}
  This is the declared model definition of action Dimension.
\end{proof}
\begin{definition}[Action Quantity]
  \label{def:physics:phyx-mini-0602:phyxminiproblems-problemphyxmini0602-actionquantity}
  \lean{PhyXMiniProblems.ProblemPhyXMini0602.ActionQuantity}
  \uses{def:physics:phyx-mini-0602:phyxminiproblems-problemphyxmini0602-actiondimension}
  A nonnegative, unit-independent action such as reduced Planck's constant
\end{definition}
\begin{proof}
  This is the declared model definition of Action Quantity.
\end{proof}
\begin{definition}[Wave Number Quantity]
  \label{def:physics:phyx-mini-0602:phyxminiproblems-problemphyxmini0602-wavenumberquantity}
  \lean{PhyXMiniProblems.ProblemPhyXMini0602.WaveNumberQuantity}
  A nonnegative wave number, with physical dimension length⁻¹
\end{definition}
\begin{proof}
  This is the declared model definition of Wave Number Quantity.
\end{proof}
\begin{definition}[delta Strength Dimension]
  \label{def:physics:phyx-mini-0602:phyxminiproblems-problemphyxmini0602-deltastrengthdimension}
  \lean{PhyXMiniProblems.ProblemPhyXMini0602.deltaStrengthDimension}
  The delta-coupling dimension energy * length = mass * length³ / time². For V(x) = -α δ(x-x₀), α has this dimension because δ has inverse length dimension
\end{definition}
\begin{proof}
  This is the declared model definition of delta Strength Dimension.
\end{proof}
\begin{definition}[Delta Strength Quantity]
  \label{def:physics:phyx-mini-0602:phyxminiproblems-problemphyxmini0602-deltastrengthquantity}
  \lean{PhyXMiniProblems.ProblemPhyXMini0602.DeltaStrengthQuantity}
  \uses{def:physics:phyx-mini-0602:phyxminiproblems-problemphyxmini0602-deltastrengthdimension}
  A nonnegative strength α of an attractive delta interaction
\end{definition}
\begin{proof}
  This is the declared model definition of Delta Strength Quantity.
\end{proof}
\begin{definition}[length In Meters]
  \label{def:physics:phyx-mini-0602:phyxminiproblems-problemphyxmini0602-lengthinmeters}
  \lean{PhyXMiniProblems.ProblemPhyXMini0602.lengthInMeters}
  \uses{def:physics:phyx-mini-0602:phyxminiproblems-problemphyxmini0602-lengthquantity}
  Metre readout of a physical length
\end{definition}
\begin{proof}
  This is the declared model definition of length In Meters.
\end{proof}
\begin{definition}[mass In Kilograms]
  \label{def:physics:phyx-mini-0602:phyxminiproblems-problemphyxmini0602-massinkilograms}
  \lean{PhyXMiniProblems.ProblemPhyXMini0602.massInKilograms}
  \uses{def:physics:phyx-mini-0602:phyxminiproblems-problemphyxmini0602-massquantity}
  Kilogram readout of a physical mass
\end{definition}
\begin{proof}
  This is the declared model definition of mass In Kilograms.
\end{proof}
\begin{definition}[energy In Joules]
  \label{def:physics:phyx-mini-0602:phyxminiproblems-problemphyxmini0602-energyinjoules}
  \lean{PhyXMiniProblems.ProblemPhyXMini0602.energyInJoules}
  Joule readout of a physical energy
\end{definition}
\begin{proof}
  This is the declared model definition of energy In Joules.
\end{proof}
\begin{definition}[action In Joule Seconds]
  \label{def:physics:phyx-mini-0602:phyxminiproblems-problemphyxmini0602-actioninjouleseconds}
  \lean{PhyXMiniProblems.ProblemPhyXMini0602.actionInJouleSeconds}
  \uses{def:physics:phyx-mini-0602:phyxminiproblems-problemphyxmini0602-actionquantity}
  Joule-second readout of an action
\end{definition}
\begin{proof}
  This is the declared model definition of action In Joule Seconds.
\end{proof}
\begin{definition}[wave Number In Inverse Meters]
  \label{def:physics:phyx-mini-0602:phyxminiproblems-problemphyxmini0602-wavenumberininversemeters}
  \lean{PhyXMiniProblems.ProblemPhyXMini0602.waveNumberInInverseMeters}
  \uses{def:physics:phyx-mini-0602:phyxminiproblems-problemphyxmini0602-wavenumberquantity}
  Inverse-metre readout of a wave number
\end{definition}
\begin{proof}
  This is the declared model definition of wave Number In Inverse Meters.
\end{proof}
\begin{definition}[delta Strength In Joule Meters]
  \label{def:physics:phyx-mini-0602:phyxminiproblems-problemphyxmini0602-deltastrengthinjoulemeters}
  \lean{PhyXMiniProblems.ProblemPhyXMini0602.deltaStrengthInJouleMeters}
  \uses{def:physics:phyx-mini-0602:phyxminiproblems-problemphyxmini0602-deltastrengthquantity}
  Joule-metre readout of the double-delta coupling α
\end{definition}
\begin{proof}
  This is the declared model definition of delta Strength In Joule Meters.
\end{proof}
\begin{definition}[Amplitude Pair]
  \label{def:physics:phyx-mini-0602:phyxminiproblems-problemphyxmini0602-amplitudepair}
  \lean{PhyXMiniProblems.ProblemPhyXMini0602.AmplitudePair}
  A column of right- and left-moving dimensionless wave amplitudes
\end{definition}
\begin{proof}
  This is the declared model definition of Amplitude Pair.
\end{proof}
\begin{definition}[Transfer Matrix]
  \label{def:physics:phyx-mini-0602:phyxminiproblems-problemphyxmini0602-transfermatrix}
  \lean{PhyXMiniProblems.ProblemPhyXMini0602.TransferMatrix}
  A complex 2 × 2 transfer matrix acting on an amplitude pair
\end{definition}
\begin{proof}
  This is the declared model definition of Transfer Matrix.
\end{proof}
\begin{definition}[Real Potential Distribution]
  \label{def:physics:phyx-mini-0602:phyxminiproblems-problemphyxmini0602-realpotentialdistribution}
  \lean{PhyXMiniProblems.ProblemPhyXMini0602.RealPotentialDistribution}
  A real scalar distribution in the metre coordinate
\end{definition}
\begin{proof}
  This is the declared model definition of Real Potential Distribution.
\end{proof}
\begin{definition}[Point Interaction]
  \label{def:physics:phyx-mini-0602:phyxminiproblems-problemphyxmini0602-pointinteraction}
  \lean{PhyXMiniProblems.ProblemPhyXMini0602.PointInteraction}
  The two point interactions, named by the M₁ and M₂ labels in the image
\end{definition}
\begin{proof}
  This is the declared model definition of Point Interaction.
\end{proof}
\begin{definition}[Free Region]
  \label{def:physics:phyx-mini-0602:phyxminiproblems-problemphyxmini0602-freeregion}
  \lean{PhyXMiniProblems.ProblemPhyXMini0602.FreeRegion}
  The three free-particle regions separated by the two point interactions
\end{definition}
\begin{proof}
  This is the declared model definition of Free Region.
\end{proof}
\begin{definition}[Supplied Double Delta Figure]
  \label{def:physics:phyx-mini-0602:phyxminiproblems-problemphyxmini0602-supplieddoubledeltafigure}
  \lean{PhyXMiniProblems.ProblemPhyXMini0602.SuppliedDoubleDeltaFigure}
  \uses{def:physics:phyx-mini-0602:phyxminiproblems-problemphyxmini0602-pointinteraction, def:physics:phyx-mini-0602:phyxminiproblems-problemphyxmini0602-freeregion}
  Raster-visible data from image 602. The image labels two separated packets M₁, M₂, places V = 0 under all three intervening regions, and labels the horizontal axis x; it does not visibly print ±a
\end{definition}
\begin{proof}
  This is the declared model definition of Supplied Double Delta Figure.
\end{proof}
\begin{definition}[Double Delta Scattering Setup]
  \label{def:physics:phyx-mini-0602:phyxminiproblems-problemphyxmini0602-doubledeltascatteringsetup}
  \lean{PhyXMiniProblems.ProblemPhyXMini0602.DoubleDeltaScatteringSetup}
  \uses{def:physics:phyx-mini-0602:phyxminiproblems-problemphyxmini0602-lengthquantity, def:physics:phyx-mini-0602:phyxminiproblems-problemphyxmini0602-massquantity, def:physics:phyx-mini-0602:phyxminiproblems-problemphyxmini0602-actionquantity, def:physics:phyx-mini-0602:phyxminiproblems-problemphyxmini0602-wavenumberquantity, def:physics:phyx-mini-0602:phyxminiproblems-problemphyxmini0602-deltastrengthquantity, def:physics:phyx-mini-0602:phyxminiproblems-problemphyxmini0602-amplitudepair, def:physics:phyx-mini-0602:phyxminiproblems-problemphyxmini0602-transfermatrix, def:physics:phyx-mini-0602:phyxminiproblems-problemphyxmini0602-realpotentialdistribution, def:physics:phyx-mini-0602:phyxminiproblems-problemphyxmini0602-pointinteraction, def:physics:phyx-mini-0602:phyxminiproblems-problemphyxmini0602-freeregion, def:physics:phyx-mini-0602:phyxminiproblems-problemphyxmini0602-supplieddoubledeltafigure}
  Independent physical quantities, amplitudes, and transfer data. The maps propagateAcrossM1 and propagateAcrossM2 act on arbitrary amplitude pairs, which makes each transfer matrix identifiable rather than constraining it on only the single displayed scattering state
\end{definition}
\begin{proof}
  This is the declared model definition of Double Delta Scattering Setup.
\end{proof}
\begin{definition}[half Separation Meters]
  \label{def:physics:phyx-mini-0602:phyxminiproblems-problemphyxmini0602-halfseparationmeters}
  \lean{PhyXMiniProblems.ProblemPhyXMini0602.halfSeparationMeters}
  \uses{def:physics:phyx-mini-0602:phyxminiproblems-problemphyxmini0602-lengthinmeters, def:physics:phyx-mini-0602:phyxminiproblems-problemphyxmini0602-doubledeltascatteringsetup}
  The half-separation a, read in metres
\end{definition}
\begin{proof}
  This is the declared model definition of half Separation Meters.
\end{proof}
\begin{definition}[particle Mass Kilograms]
  \label{def:physics:phyx-mini-0602:phyxminiproblems-problemphyxmini0602-particlemasskilograms}
  \lean{PhyXMiniProblems.ProblemPhyXMini0602.particleMassKilograms}
  \uses{def:physics:phyx-mini-0602:phyxminiproblems-problemphyxmini0602-massinkilograms, def:physics:phyx-mini-0602:phyxminiproblems-problemphyxmini0602-doubledeltascatteringsetup}
  The particle mass m, read in kilograms
\end{definition}
\begin{proof}
  This is the declared model definition of particle Mass Kilograms.
\end{proof}
\begin{definition}[delta Strength SI]
  \label{def:physics:phyx-mini-0602:phyxminiproblems-problemphyxmini0602-deltastrengthsi}
  \lean{PhyXMiniProblems.ProblemPhyXMini0602.deltaStrengthSI}
  \uses{def:physics:phyx-mini-0602:phyxminiproblems-problemphyxmini0602-deltastrengthinjoulemeters, def:physics:phyx-mini-0602:phyxminiproblems-problemphyxmini0602-doubledeltascatteringsetup}
  The attractive coupling α, read in joule-metres
\end{definition}
\begin{proof}
  This is the declared model definition of delta Strength SI.
\end{proof}
\begin{definition}[reduced Planck SI]
  \label{def:physics:phyx-mini-0602:phyxminiproblems-problemphyxmini0602-reducedplancksi}
  \lean{PhyXMiniProblems.ProblemPhyXMini0602.reducedPlanckSI}
  \uses{def:physics:phyx-mini-0602:phyxminiproblems-problemphyxmini0602-actioninjouleseconds, def:physics:phyx-mini-0602:phyxminiproblems-problemphyxmini0602-doubledeltascatteringsetup}
  The reduced Planck action ℏ, read in joule-seconds
\end{definition}
\begin{proof}
  This is the declared model definition of reduced Planck SI.
\end{proof}
\begin{definition}[incident Energy Joules]
  \label{def:physics:phyx-mini-0602:phyxminiproblems-problemphyxmini0602-incidentenergyjoules}
  \lean{PhyXMiniProblems.ProblemPhyXMini0602.incidentEnergyJoules}
  \uses{def:physics:phyx-mini-0602:phyxminiproblems-problemphyxmini0602-energyinjoules, def:physics:phyx-mini-0602:phyxminiproblems-problemphyxmini0602-doubledeltascatteringsetup}
  The incident energy E, read in joules
\end{definition}
\begin{proof}
  This is the declared model definition of incident Energy Joules.
\end{proof}
\begin{definition}[wave Number SI]
  \label{def:physics:phyx-mini-0602:phyxminiproblems-problemphyxmini0602-wavenumbersi}
  \lean{PhyXMiniProblems.ProblemPhyXMini0602.waveNumberSI}
  \uses{def:physics:phyx-mini-0602:phyxminiproblems-problemphyxmini0602-wavenumberininversemeters, def:physics:phyx-mini-0602:phyxminiproblems-problemphyxmini0602-doubledeltascatteringsetup}
  The wave number k, read in inverse metres
\end{definition}
\begin{proof}
  This is the declared model definition of wave Number SI.
\end{proof}
\begin{definition}[amplitude A]
  \label{def:physics:phyx-mini-0602:phyxminiproblems-problemphyxmini0602-amplitudea}
  \lean{PhyXMiniProblems.ProblemPhyXMini0602.amplitudeA}
  \uses{def:physics:phyx-mini-0602:phyxminiproblems-problemphyxmini0602-doubledeltascatteringsetup}
  Source amplitude A, incident from the left
\end{definition}
\begin{proof}
  This is the declared model definition of amplitude A.
\end{proof}
\begin{definition}[amplitude B]
  \label{def:physics:phyx-mini-0602:phyxminiproblems-problemphyxmini0602-amplitudeb}
  \lean{PhyXMiniProblems.ProblemPhyXMini0602.amplitudeB}
  \uses{def:physics:phyx-mini-0602:phyxminiproblems-problemphyxmini0602-doubledeltascatteringsetup}
  Source amplitude B, outgoing to the left
\end{definition}
\begin{proof}
  This is the declared model definition of amplitude B.
\end{proof}
\begin{definition}[amplitude F]
  \label{def:physics:phyx-mini-0602:phyxminiproblems-problemphyxmini0602-amplitudef}
  \lean{PhyXMiniProblems.ProblemPhyXMini0602.amplitudeF}
  \uses{def:physics:phyx-mini-0602:phyxminiproblems-problemphyxmini0602-doubledeltascatteringsetup}
  Source amplitude F, outgoing to the right
\end{definition}
\begin{proof}
  This is the declared model definition of amplitude F.
\end{proof}
\begin{definition}[amplitude G]
  \label{def:physics:phyx-mini-0602:phyxminiproblems-problemphyxmini0602-amplitudeg}
  \lean{PhyXMiniProblems.ProblemPhyXMini0602.amplitudeG}
  \uses{def:physics:phyx-mini-0602:phyxminiproblems-problemphyxmini0602-doubledeltascatteringsetup}
  Source amplitude G, incoming from the right
\end{definition}
\begin{proof}
  This is the declared model definition of amplitude G.
\end{proof}
\begin{definition}[Matches Double Delta Scenario]
  \label{def:physics:phyx-mini-0602:phyxminiproblems-problemphyxmini0602-matchesdoubledeltascenario}
  \lean{PhyXMiniProblems.ProblemPhyXMini0602.MatchesDoubleDeltaScenario}
  \uses{def:physics:phyx-mini-0602:phyxminiproblems-problemphyxmini0602-doubledeltascatteringsetup, def:physics:phyx-mini-0602:phyxminiproblems-problemphyxmini0602-halfseparationmeters, def:physics:phyx-mini-0602:phyxminiproblems-problemphyxmini0602-deltastrengthsi}
  The potential stated in the prose, including both locations and every zero-potential region. The distribution is a coherent-SI representation of -α[δ(x+a)+δ(x-a)]; no pointwise value is assigned at a delta singularity
\end{definition}
\begin{proof}
  This is the declared model definition of Matches Double Delta Scenario.
\end{proof}
\begin{definition}[Matches Supplied Double Delta Figure]
  \label{def:physics:phyx-mini-0602:phyxminiproblems-problemphyxmini0602-matchessupplieddoubledeltafigure}
  \lean{PhyXMiniProblems.ProblemPhyXMini0602.MatchesSuppliedDoubleDeltaFigure}
  \uses{def:physics:phyx-mini-0602:phyxminiproblems-problemphyxmini0602-doubledeltascatteringsetup}
  All labels and qualitative features visible in the supplied raster
\end{definition}
\begin{proof}
  This is the declared model definition of Matches Supplied Double Delta Figure.
\end{proof}
\begin{definition}[Matches Left Incident Scattering State]
  \label{def:physics:phyx-mini-0602:phyxminiproblems-problemphyxmini0602-matchesleftincidentscatteringstate}
  \lean{PhyXMiniProblems.ProblemPhyXMini0602.MatchesLeftIncidentScatteringState}
  \uses{def:physics:phyx-mini-0602:phyxminiproblems-problemphyxmini0602-doubledeltascatteringsetup, def:physics:phyx-mini-0602:phyxminiproblems-problemphyxmini0602-amplitudea, def:physics:phyx-mini-0602:phyxminiproblems-problemphyxmini0602-amplitudeg}
  Incidence is from the left: A ≠ 0 and there is no incoming right wave G
\end{definition}
\begin{proof}
  This is the declared model definition of Matches Left Incident Scattering State.
\end{proof}
\begin{definition}[Has Physical Double Delta Parameters]
  \label{def:physics:phyx-mini-0602:phyxminiproblems-problemphyxmini0602-hasphysicaldoubledeltaparameters}
  \lean{PhyXMiniProblems.ProblemPhyXMini0602.HasPhysicalDoubleDeltaParameters}
  \uses{def:physics:phyx-mini-0602:phyxminiproblems-problemphyxmini0602-doubledeltascatteringsetup, def:physics:phyx-mini-0602:phyxminiproblems-problemphyxmini0602-halfseparationmeters, def:physics:phyx-mini-0602:phyxminiproblems-problemphyxmini0602-particlemasskilograms, def:physics:phyx-mini-0602:phyxminiproblems-problemphyxmini0602-deltastrengthsi, def:physics:phyx-mini-0602:phyxminiproblems-problemphyxmini0602-reducedplancksi, def:physics:phyx-mini-0602:phyxminiproblems-problemphyxmini0602-incidentenergyjoules, def:physics:phyx-mini-0602:phyxminiproblems-problemphyxmini0602-wavenumbersi}
  Positivity and probability-range conditions on the physical parameters
\end{definition}
\begin{proof}
  This is the declared model definition of Has Physical Double Delta Parameters.
\end{proof}
\begin{definition}[Uses Standard Reduced Planck Constant]
  \label{def:physics:phyx-mini-0602:phyxminiproblems-problemphyxmini0602-usesstandardreducedplanckconstant}
  \lean{PhyXMiniProblems.ProblemPhyXMini0602.UsesStandardReducedPlanckConstant}
  \uses{def:physics:phyx-mini-0602:phyxminiproblems-problemphyxmini0602-doubledeltascatteringsetup, def:physics:phyx-mini-0602:phyxminiproblems-problemphyxmini0602-reducedplancksi}
  Physlib's Constants.ℏ is the standard reduced Planck constant, expressed in joule-seconds. This calibration does not constrain a requested matrix entry or transmission formula
\end{definition}
\begin{proof}
  This is the declared model definition of Uses Standard Reduced Planck Constant.
\end{proof}
\begin{definition}[plane Wave Value At]
  \label{def:physics:phyx-mini-0602:phyxminiproblems-problemphyxmini0602-planewavevalueat}
  \lean{PhyXMiniProblems.ProblemPhyXMini0602.planeWaveValueAt}
  \uses{def:physics:phyx-mini-0602:phyxminiproblems-problemphyxmini0602-amplitudepair}
  The value of a two-component plane-wave expansion at a metre coordinate
\end{definition}
\begin{proof}
  This is the declared model definition of plane Wave Value At.
\end{proof}
\begin{definition}[plane Wave Derivative At]
  \label{def:physics:phyx-mini-0602:phyxminiproblems-problemphyxmini0602-planewavederivativeat}
  \lean{PhyXMiniProblems.ProblemPhyXMini0602.planeWaveDerivativeAt}
  \uses{def:physics:phyx-mini-0602:phyxminiproblems-problemphyxmini0602-amplitudepair}
  The spatial derivative of the same plane-wave expansion
\end{definition}
\begin{proof}
  This is the declared model definition of plane Wave Derivative At.
\end{proof}
\begin{definition}[Satisfies Attractive Delta Matching]
  \label{def:physics:phyx-mini-0602:phyxminiproblems-problemphyxmini0602-satisfiesattractivedeltamatching}
  \lean{PhyXMiniProblems.ProblemPhyXMini0602.SatisfiesAttractiveDeltaMatching}
  \uses{def:physics:phyx-mini-0602:phyxminiproblems-problemphyxmini0602-amplitudepair, def:physics:phyx-mini-0602:phyxminiproblems-problemphyxmini0602-doubledeltascatteringsetup, def:physics:phyx-mini-0602:phyxminiproblems-problemphyxmini0602-particlemasskilograms, def:physics:phyx-mini-0602:phyxminiproblems-problemphyxmini0602-deltastrengthsi, def:physics:phyx-mini-0602:phyxminiproblems-problemphyxmini0602-reducedplancksi, def:physics:phyx-mini-0602:phyxminiproblems-problemphyxmini0602-wavenumbersi, def:physics:phyx-mini-0602:phyxminiproblems-problemphyxmini0602-planewavevalueat, def:physics:phyx-mini-0602:phyxminiproblems-problemphyxmini0602-planewavederivativeat}
  Continuity and the derivative jump at one attractive point interaction. For V = -α δ(x-x₀), the jump is ψ'(x₀⁺)-ψ'(x₀⁻) = -(2mα/ℏ²) ψ(x₀). This general boundary law contains no candidate transfer-matrix formula
\end{definition}
\begin{proof}
  This is the declared model definition of Satisfies Attractive Delta Matching.
\end{proof}
\begin{definition}[Satisfies Free Particle Dispersion Law]
  \label{def:physics:phyx-mini-0602:phyxminiproblems-problemphyxmini0602-satisfiesfreeparticledispersionlaw}
  \lean{PhyXMiniProblems.ProblemPhyXMini0602.SatisfiesFreeParticleDispersionLaw}
  \uses{def:physics:phyx-mini-0602:phyxminiproblems-problemphyxmini0602-doubledeltascatteringsetup, def:physics:phyx-mini-0602:phyxminiproblems-problemphyxmini0602-particlemasskilograms, def:physics:phyx-mini-0602:phyxminiproblems-problemphyxmini0602-reducedplancksi, def:physics:phyx-mini-0602:phyxminiproblems-problemphyxmini0602-incidentenergyjoules, def:physics:phyx-mini-0602:phyxminiproblems-problemphyxmini0602-wavenumbersi}
  The nonrelativistic free-particle dispersion relation ℏ²k² = 2mE
\end{definition}
\begin{proof}
  This is the declared model definition of Satisfies Free Particle Dispersion Law.
\end{proof}
\begin{definition}[Satisfies Double Delta Scattering Laws]
  \label{def:physics:phyx-mini-0602:phyxminiproblems-problemphyxmini0602-satisfiesdoubledeltascatteringlaws}
  \lean{PhyXMiniProblems.ProblemPhyXMini0602.SatisfiesDoubleDeltaScatteringLaws}
  \uses{def:physics:phyx-mini-0602:phyxminiproblems-problemphyxmini0602-doubledeltascatteringsetup, def:physics:phyx-mini-0602:phyxminiproblems-problemphyxmini0602-amplitudea, def:physics:phyx-mini-0602:phyxminiproblems-problemphyxmini0602-amplitudef, def:physics:phyx-mini-0602:phyxminiproblems-problemphyxmini0602-satisfiesattractivedeltamatching}
  The two matching operations act on arbitrary amplitudes, their matrices implement those operations, and the total matrix implements their sequential action. The final two fields specialize this generic propagation to the named amplitudes and define transmission as transmitted-to-incident flux for equal asymptotic wave numbers. No matrix entry or closed transmission answer is assumed here
\end{definition}
\begin{proof}
  This is the declared model definition of Satisfies Double Delta Scattering Laws.
\end{proof}
\begin{definition}[dimensionless Delta Coupling]
  \label{def:physics:phyx-mini-0602:phyxminiproblems-problemphyxmini0602-dimensionlessdeltacoupling}
  \lean{PhyXMiniProblems.ProblemPhyXMini0602.dimensionlessDeltaCoupling}
  \uses{def:physics:phyx-mini-0602:phyxminiproblems-problemphyxmini0602-doubledeltascatteringsetup, def:physics:phyx-mini-0602:phyxminiproblems-problemphyxmini0602-particlemasskilograms, def:physics:phyx-mini-0602:phyxminiproblems-problemphyxmini0602-deltastrengthsi, def:physics:phyx-mini-0602:phyxminiproblems-problemphyxmini0602-reducedplancksi, def:physics:phyx-mini-0602:phyxminiproblems-problemphyxmini0602-wavenumbersi}
  The dimensionless attractive-delta coupling η = mα/(ℏ²k)
\end{definition}
\begin{proof}
  This is the declared model definition of dimensionless Delta Coupling.
\end{proof}
\begin{definition}[attractive Delta Transfer Matrix]
  \label{def:physics:phyx-mini-0602:phyxminiproblems-problemphyxmini0602-attractivedeltatransfermatrix}
  \lean{PhyXMiniProblems.ProblemPhyXMini0602.attractiveDeltaTransferMatrix}
  \uses{def:physics:phyx-mini-0602:phyxminiproblems-problemphyxmini0602-transfermatrix, def:physics:phyx-mini-0602:phyxminiproblems-problemphyxmini0602-doubledeltascatteringsetup, def:physics:phyx-mini-0602:phyxminiproblems-problemphyxmini0602-wavenumbersi, def:physics:phyx-mini-0602:phyxminiproblems-problemphyxmini0602-dimensionlessdeltacoupling}
  Candidate transfer matrix for one attractive delta at positionMeters. This is an answer expression, not a field of any premise structure
\end{definition}
\begin{proof}
  This is the declared model definition of attractive Delta Transfer Matrix.
\end{proof}
\begin{definition}[explicit Double Delta Transfer Matrix]
  \label{def:physics:phyx-mini-0602:phyxminiproblems-problemphyxmini0602-explicitdoubledeltatransfermatrix}
  \lean{PhyXMiniProblems.ProblemPhyXMini0602.explicitDoubleDeltaTransferMatrix}
  \uses{def:physics:phyx-mini-0602:phyxminiproblems-problemphyxmini0602-transfermatrix, def:physics:phyx-mini-0602:phyxminiproblems-problemphyxmini0602-doubledeltascatteringsetup, def:physics:phyx-mini-0602:phyxminiproblems-problemphyxmini0602-halfseparationmeters, def:physics:phyx-mini-0602:phyxminiproblems-problemphyxmini0602-wavenumbersi, def:physics:phyx-mini-0602:phyxminiproblems-problemphyxmini0602-dimensionlessdeltacoupling}
  The explicit total transfer matrix obtained in the source convention (F,G)ᵀ = M(A,B)ᵀ. Lean indices (0,0), (0,1), (1,0), (1,1) correspond to the source's M₁₁, M₁₂, M₂₁, M₂₂
\end{definition}
\begin{proof}
  This is the declared model definition of explicit Double Delta Transfer Matrix.
\end{proof}
\begin{definition}[Answer Choice]
  \label{def:physics:phyx-mini-0602:phyxminiproblems-problemphyxmini0602-answerchoice}
  \lean{PhyXMiniProblems.ProblemPhyXMini0602.AnswerChoice}
  Labels of the four expressions printed in the source
\end{definition}
\begin{proof}
  This is the declared model definition of Answer Choice.
\end{proof}
\begin{definition}[displayed Transmission Coefficient]
  \label{def:physics:phyx-mini-0602:phyxminiproblems-problemphyxmini0602-displayedtransmissioncoefficient}
  \lean{PhyXMiniProblems.ProblemPhyXMini0602.displayedTransmissionCoefficient}
  \uses{def:physics:phyx-mini-0602:phyxminiproblems-problemphyxmini0602-doubledeltascatteringsetup, def:physics:phyx-mini-0602:phyxminiproblems-problemphyxmini0602-halfseparationmeters, def:physics:phyx-mini-0602:phyxminiproblems-problemphyxmini0602-particlemasskilograms, def:physics:phyx-mini-0602:phyxminiproblems-problemphyxmini0602-deltastrengthsi, def:physics:phyx-mini-0602:phyxminiproblems-problemphyxmini0602-reducedplancksi, def:physics:phyx-mini-0602:phyxminiproblems-problemphyxmini0602-incidentenergyjoules, def:physics:phyx-mini-0602:phyxminiproblems-problemphyxmini0602-wavenumbersi, def:physics:phyx-mini-0602:phyxminiproblems-problemphyxmini0602-answerchoice}
  The literal coefficient printed beside each answer label. These are scalar readout formulas; no option is used as a premise
\end{definition}
\begin{proof}
  This is the declared model definition of displayed Transmission Coefficient.
\end{proof}
\begin{definition}[recorded Dataset Answer]
  \label{def:physics:phyx-mini-0602:phyxminiproblems-problemphyxmini0602-recordeddatasetanswer}
  \lean{PhyXMiniProblems.ProblemPhyXMini0602.recordedDatasetAnswer}
  \uses{def:physics:phyx-mini-0602:phyxminiproblems-problemphyxmini0602-answerchoice}
  Answer label recorded by the dataset; this is metadata, not a premise
\end{definition}
\begin{proof}
  This is the declared model definition of recorded Dataset Answer.
\end{proof}
\begin{definition}[physically Supported Transmission Coefficient]
  \label{def:physics:phyx-mini-0602:phyxminiproblems-problemphyxmini0602-physicallysupportedtransmissioncoefficient}
  \lean{PhyXMiniProblems.ProblemPhyXMini0602.physicallySupportedTransmissionCoefficient}
  \uses{def:physics:phyx-mini-0602:phyxminiproblems-problemphyxmini0602-doubledeltascatteringsetup, def:physics:phyx-mini-0602:phyxminiproblems-problemphyxmini0602-halfseparationmeters, def:physics:phyx-mini-0602:phyxminiproblems-problemphyxmini0602-particlemasskilograms, def:physics:phyx-mini-0602:phyxminiproblems-problemphyxmini0602-deltastrengthsi, def:physics:phyx-mini-0602:phyxminiproblems-problemphyxmini0602-reducedplancksi, def:physics:phyx-mini-0602:phyxminiproblems-problemphyxmini0602-incidentenergyjoules, def:physics:phyx-mini-0602:phyxminiproblems-problemphyxmini0602-wavenumbersi}
  The transmission coefficient supported by the delta matching law. With η = mα/(ℏ²k) and ka = k a, this is T = [1 + 4η² (cos(2ka) - η sin(2ka))²]⁻¹. Using ℏ²k² = 2mE, the prefactor 4η² is equivalently 2mα²/(ℏ²E). Unlike the literal source option C, both the inner coupling and the squared trigonometric factor are dimensionally and algebraically consistent with T = 1 / |M₂₂|²
\end{definition}
\begin{proof}
  This is the declared model definition of physically Supported Transmission Coefficient.
\end{proof}
% --- Archon declaration topology end ---
% --- Archon physics formalization source end ---
